\documentclass[aspectratio=169,10pt]{beamer}
% Shared Beamer setup for Fall 2026 course slides.
% Clean Cal Poly-inspired theme: no custom class files or uncommon packages.
\mode<presentation>{
  \usetheme{default}
  \usefonttheme{professionalfonts}
}

\definecolor{CalPolyGreen}{HTML}{154734}
\definecolor{CalPolyGold}{HTML}{BD8B13}
\definecolor{CalPolySage}{HTML}{7FA28F}
\definecolor{CalPolyMint}{HTML}{DDF1E7}
\definecolor{CalPolySand}{HTML}{A29F86}
\definecolor{CalPolyBeige}{HTML}{EFEEDE}
\definecolor{DeckInk}{HTML}{1F2933}
\definecolor{DeckMuted}{HTML}{5F6B7A}
\definecolor{DeckSoft}{HTML}{F7F7F5}

\setbeamersize{text margin left=0.55in,text margin right=0.55in}
\setbeamercolor{normal text}{fg=DeckInk,bg=white}
\setbeamercolor{structure}{fg=CalPolyGreen}
\setbeamercolor{title}{fg=white}
\setbeamercolor{frametitle}{fg=CalPolyGreen,bg=white}
\setbeamercolor{block title}{bg=CalPolySage,fg=white}
\setbeamercolor{block body}{bg=CalPolyMint,fg=black}
\setbeamercolor{alerted text}{fg=CalPolyGold}

\setbeamerfont{title}{size=\Huge,series=\bfseries}
\setbeamerfont{frametitle}{size=\Large,series=\bfseries}
\setbeamerfont{block title}{size=\normalsize,series=\bfseries}

\setbeamertemplate{navigation symbols}{}
\setbeamertemplate{itemize item}{\color{CalPolyGold}$\blacktriangleright$}
\setbeamertemplate{itemize subitem}{\color{CalPolyGreen}$\bullet$}
\setbeamertemplate{footline}{
  \hfill{\color{DeckMuted}\scriptsize\insertframenumber/\inserttotalframenumber}
  \hspace{0.18in}\vspace{0.12in}
}
\setbeamertemplate{frametitle}{
  \vspace{0.18in}
  \hspace{0.02in}\insertframetitle\par
  \vspace{0.08in}
  \noindent\makebox[\linewidth][l]{%
    {\color{CalPolyGold}\rule{0.55in}{2pt}}%
    {\color{CalPolyGreen}\rule{\dimexpr\linewidth-0.55in\relax}{2pt}}%
  }
}

\newcommand{\CourseNumber}{Course}
\newcommand{\CourseTitle}{Course Title}
\newcommand{\LectureNumber}{0}
\newcommand{\LectureTitle}{Lecture Title}
\newcommand{\LectureDate}{Date}
\newcommand{\CourseUnit}{Unit}

\newcommand{\coursenumber}[1]{\renewcommand{\CourseNumber}{#1}}
\newcommand{\coursetitle}[1]{\renewcommand{\CourseTitle}{#1}}
\newcommand{\lecturenumber}[1]{\renewcommand{\LectureNumber}{#1}}
\newcommand{\lecturetitle}[1]{\renewcommand{\LectureTitle}{#1}}
\newcommand{\lecturedate}[1]{\renewcommand{\LectureDate}{#1}}
\newcommand{\courseunit}[1]{\renewcommand{\CourseUnit}{#1}}

\author{Kirk Duran}
\institute{California Polytechnic State University, San Luis Obispo}
\date{\LectureDate}

\newcommand{\makedecktitle}{
  \begingroup
  \setbeamercolor{background canvas}{bg=CalPolyGreen}
  \begin{frame}[plain]
    \vfill
    \begin{center}
      {\usebeamerfont{title}\color{white}\LectureTitle\par}
      \vspace{0.18in}
      {\Large\color{white}Lecture \LectureNumber\par}
    \end{center}
    \vfill
    \noindent{\color{white}\rule{\linewidth}{1.2pt}}\par
    \vspace{0.08in}
    \noindent\colorbox{CalPolyGold}{\makebox[0.24\linewidth][c]{\color{white}\Large\CourseNumber}}
    \hspace{0.12in}{\color{white}\Large Kirk Duran}
  \end{frame}
  \endgroup
}

\newenvironment{keyidea}{\begin{block}{Key idea}}{\end{block}}
\newenvironment{activity}{\begin{block}{In class}}{\end{block}}
\newenvironment{cpdefinition}{\begin{block}{Definition}}{\end{block}}
\newenvironment{exampleblockcp}{
  \begingroup
  \setbeamercolor{block title}{bg=CalPolySand,fg=white}
  \setbeamercolor{block body}{bg=CalPolyBeige,fg=black}
  \begin{block}{Example}
}{
  \end{block}
  \endgroup
}

\newcommand{\imageplaceholder}[2][0.3in]{%
  \noindent\makebox[\linewidth][c]{%
    \fcolorbox{CalPolySage}{CalPolyBeige}{%
      \parbox[c][#1][c]{0.86\linewidth}{%
        \centering
        {\color{DeckMuted}\scriptsize Image placeholder: }%
        {\color{DeckInk}\scriptsize\ttfamily\detokenize{#2}}%
      }%
    }%
  }%
}

\newcommand{\sectionframe}[1]{
  \begingroup
  \setbeamercolor{background canvas}{bg=CalPolyGreen}
  \begin{frame}[plain]
    \vfill
    {\color{white}\LARGE\bfseries #1\par}
    \vspace{0.18in}
    {\color{CalPolyGold}\rule{0.22\paperwidth}{3pt}}
    {\color{white}\rule{0.62\paperwidth}{3pt}}
    \vfill
  \end{frame}
  \endgroup
}

\newcommand{\placeholdernote}{
  \begin{block}{Placeholder}
    Replace these bullets with the final lecture material, examples, and in-class activities.
  \end{block}
}

\usepackage[T1]{fontenc}
\usepackage{lmodern}
\pdfstringdefDisableCommands{\def\translate#1{#1}}
\graphicspath{{../assets/lecture-31/}}
\setbeamersize{text margin left=.42in,text margin right=.42in}

% Adapted from the instructor's CS1001_Lecture_31_Sorting_Algorithms.pdf.
% All 14 supplied array images are preserved; trim removes only their white margins.
% Green: sorted prefix. Gray: unprocessed. Gold: active value or insertion gap.
\newcommand{\arraydiagram}[2][1.65in]{%
  \raisebox{-.43\height}{\includegraphics[trim=6.3bp 6.3bp 6.3bp 6.3bp,clip,width=#1]{array-#2.png}}%
}
\coursenumber{CS 1001}
\coursetitle{Introduction to Computing}
\lecturenumber{31}
\lecturetitle{Sorting Algorithms}
\lecturedate{November 9, 2026}
\courseunit{12. Algorithms}

\begin{document}
\makedecktitle

\begin{frame}{Today}
\small
\begin{itemize}
 \item Connect binary search from last lecture to the need for ordered data.
 \item Understand sorting as the growth of a sorted region.
 \item Perform selection sort by hand, implement it, and trace the Python code.
 \item Perform insertion sort by hand, implement it, and trace the Python code.
 \item Use Python's \texttt{.sort()} and \texttt{sorted()} for ordinary programs.
 \item Explain \texttt{lambda} and use \texttt{key=} to sort lists of objects.
\end{itemize}
\begin{keyidea}First learn how sorting works; then use Python's built-in tools for ordinary programs.\end{keyidea}
\end{frame}

\begin{frame}{From Last Time: Binary Search Needs Sorted Data}
\small
\begin{itemize}
 \item Of the list-search algorithms we studied, binary search uses fewer checks on large lists when the list is already ordered.
 \item Binary search can decide which half to discard only because values are in order.
 \item What happens when the collection we receive is not sorted?
\end{itemize}
\begin{center}\arraydiagram{01}\quad$\longrightarrow$\quad\arraydiagram{02}\end{center}
\begin{keyidea}Before we can use binary search on an unordered list, we need to put the list in order. Sorting itself also takes work.\end{keyidea}
\end{frame}

\sectionframe{Part 1: The Idea of Sorting}

\begin{frame}{What Is Sorting?}
\small
\begin{cpdefinition}Sorting arranges elements according to an ordering rule, such as increasing numeric value or alphabetical name.\end{cpdefinition}
\begin{itemize}
 \item Numbers can be arranged from smallest to largest.
 \item Names can be arranged in alphabetical order.
 \item Objects can be arranged by a selected attribute, such as \texttt{student.gpa}.
\end{itemize}
\begin{block}{Our convention today}Both algorithms arrange a list of integers in ascending order and modify that list directly.\end{block}
\end{frame}

\begin{frame}{Visualizing a Growing Sorted Region}
\small
\begin{center}
\begin{tabular}{ll}
Initial list & \arraydiagram{01}\\[7pt]
Some work completed & \arraydiagram{03}\\[7pt]
More work completed & \arraydiagram{04}\\[7pt]
Finished & \arraydiagram{02}
\end{tabular}
\end{center}
\begin{itemize}
 \item Green marks the sorted prefix; gray marks the unprocessed region.
 \item Gold will mark the active value or gap in a trace.
\end{itemize}
\begin{keyidea}Selection sort and insertion sort both grow a sorted prefix, but grow it differently.\end{keyidea}
\end{frame}

\begin{frame}{Two Ways to Grow the Sorted Region}
\small
\begin{columns}[T,onlytextwidth]
\begin{column}{.48\textwidth}
\begin{block}{Selection sort}Find the smallest remaining value and move it to the next position in the sorted region.\end{block}
\begin{center}\arraydiagram{05}\end{center}
\end{column}
\begin{column}{.48\textwidth}
\begin{block}{Insertion sort}Take the next unsorted value and insert it into the correct position within the sorted region.\end{block}
\begin{center}\arraydiagram{06}\end{center}
\end{column}
\end{columns}
\begin{keyidea}Selection sort's prefix contains the smallest finished values. Insertion sort's prefix is ordered, but later values may be smaller.\end{keyidea}
\end{frame}

\sectionframe{Part 2: Selection Sort}

\begin{frame}{Selection Sort: Example Before Code}
\small
Start with \texttt{[5, 2, 4, 1, 3]}. Find the smallest value in the whole unsorted region.
\begin{center}\arraydiagram[1.9in]{07}\end{center}
\begin{itemize}
 \item The smallest value is \texttt{1} at index \texttt{3}.
 \item Swap that value with the first unsorted position (index \texttt{0}).
\end{itemize}
\begin{center}\arraydiagram[1.9in]{05}\end{center}
\begin{keyidea}One complete pass puts one more value in its final sorted position.\end{keyidea}
\end{frame}

\begin{frame}{Selection Sort: Continue the Example}
\small
\begin{center}
\begin{tabular}{lll}
Beginning & \arraydiagram[1.5in]{01} & unsorted\\[6pt]
Pass 1: choose 1 & \arraydiagram[1.5in]{05} & first position fixed\\[6pt]
Pass 2: choose 2 & \arraydiagram[1.5in]{03} & same-position exchange\\[6pt]
Pass 3: choose 3 & \arraydiagram[1.5in]{04} & swap 4 and 3\\[6pt]
Pass 4: choose 4 & \arraydiagram[1.5in]{02} & finished
\end{tabular}
\end{center}
\begin{block}{Observation}We never need to examine the finished green positions when searching for the next minimum.\end{block}
\end{frame}

\begin{frame}{Selection Sort: Describe the Procedure}
\small
\begin{enumerate}
 \item Mark the first unsorted position as \texttt{start}.
 \item Assume the smallest remaining value is at \texttt{start}.
 \item Scan the remaining unsorted positions to find the actual minimum.
 \item Swap that minimum into position \texttt{start}.
 \item Move \texttt{start} one position to the right and repeat.
\end{enumerate}
\begin{keyidea}\textbf{Selection:} repeatedly select the minimum from the remaining values.\end{keyidea}
\end{frame}

\begin{frame}[fragile]{Selection Sort: Python Implementation}
\small
\begin{block}{Sort the existing list in ascending order}
\footnotesize
\begin{verbatim}
def selection_sort(values: list[int]) -> None:
    for start in range(len(values) - 1):
        min_index = start
        for scan in range(start + 1, len(values)):
            if values[scan] < values[min_index]:
                min_index = scan
        values[start], values[min_index] = values[min_index], values[start]
\end{verbatim}
\end{block}
\begin{itemize}
 \item \texttt{start}: the position we are about to finish.
 \item \texttt{min\_index}: the smallest value's index seen during this pass.
 \item \texttt{scan}: checks the remaining candidates.
 \item The assignment exchanges the two values after the scan ends.
\end{itemize}
\end{frame}

\begin{frame}[fragile]{Trace the Code: Selection Sort, First Pass}
\small
\begin{columns}[T,onlytextwidth]
\begin{column}{.49\textwidth}
\begin{block}{Relevant code}
\footnotesize
\begin{verbatim}
start = 0
min_index = start
for scan in range(1, 5):
    if values[scan] < values[min_index]:
        min_index = scan
\end{verbatim}
\end{block}
\begin{center}\arraydiagram{07}\end{center}
\end{column}
\begin{column}{.48\textwidth}
\begin{block}{Variable trace}
\begin{tabular}{ccl}
\texttt{scan} & \texttt{min\_index} & Why?\\\hline
1 & 1 & $2<5$\\
2 & 1 & $4<2$ false\\
3 & 3 & $1<2$\\
4 & 3 & $3<1$ false
\end{tabular}
\end{block}
\end{column}
\end{columns}
\begin{keyidea}The scan remembers the index of the minimum; it does not swap on every comparison.\end{keyidea}
\end{frame}

\begin{frame}[fragile]{Trace the Code: Selection Sort, The Swap}
\small
\begin{block}{The line that runs after the inner loop}
\footnotesize
\begin{verbatim}
values[start], values[min_index] = values[min_index], values[start]
\end{verbatim}
\end{block}
\begin{center}\arraydiagram{07}\quad$\longrightarrow$\quad\arraydiagram{05}\end{center}
\begin{itemize}
 \item At this moment, \texttt{start = 0} and \texttt{min\_index = 3}.
 \item The values at positions \texttt{0} and \texttt{3} exchange places.
 \item Python evaluates both right-hand values before assigning the left-hand positions.
 \item The next iteration begins with \texttt{start = 1}.
\end{itemize}
\end{frame}

\begin{frame}{Trace the Code: Selection Sort, Outer Loop}
\small
\begin{center}
\begin{tabular}{ccll}
\texttt{start} & \texttt{min\_index} & Action & List after pass\\\hline
0 & 3 & swap 5,1 & \texttt{[1,2,4,5,3]}\\
1 & 1 & same position & \texttt{[1,2,4,5,3]}\\
2 & 4 & swap 4,3 & \texttt{[1,2,3,5,4]}\\
3 & 4 & swap 5,4 & \texttt{[1,2,3,4,5]}
\end{tabular}
\end{center}
\begin{keyidea}\texttt{start} also counts the positions already finished. When only one value remains, it must already be in place.\end{keyidea}
\end{frame}

\begin{frame}{Selection Sort: Check Understanding}
\small
\begin{center}\arraydiagram[2in]{03}\end{center}
\begin{enumerate}
 \item Which positions are already finished?
 \item What value should be selected next from the remaining region?
 \item Which two indices should be swapped?
 \item What will the list look like after the next pass?
\end{enumerate}
\begin{alertblock}{Remember}Choose the minimum from the unsorted suffix, not from the whole list again.\end{alertblock}
\end{frame}

\sectionframe{Part 3: Insertion Sort}

\begin{frame}{Insertion Sort: The Playing-Card Idea}
\small
\begin{itemize}
 \item Imagine holding cards in a hand that is already sorted.
 \item Pick up one more card and slide larger cards to the right to make room.
 \item Place the new card into the open spot.
\end{itemize}
\begin{center}\arraydiagram{08}\quad$\longrightarrow$\quad\arraydiagram{06}\end{center}
\begin{keyidea}\textbf{Insertion:} take the next value and insert it into the correct position within the sorted prefix.\end{keyidea}
\end{frame}

\begin{frame}{Insertion Sort: Example Before Code}
\small
Start with \texttt{[5, 2, 4, 1, 3]}. A one-element prefix is already sorted.
\begin{center}
\begin{tabular}{ll}
One sorted value & \arraydiagram[1.5in]{08}\\[5pt]
Insert 2 before 5 & \arraydiagram[1.5in]{06}\\[5pt]
Insert 4 between 2, 5 & \arraydiagram[1.5in]{09}\\[5pt]
Insert 1 at the beginning & \arraydiagram[1.5in]{10}\\[5pt]
Insert 3 in the middle & \arraydiagram[1.5in]{02}
\end{tabular}
\end{center}
\begin{block}{Notice}The prefix is ordered after every insertion, but need not contain the smallest values in the entire list yet.\end{block}
\end{frame}

\begin{frame}{Insertion Sort: What Does ``Insert'' Mean?}
\small
Suppose the list is \texttt{[2, 4, 5, 1, 3]} and the saved key is \texttt{1}.
\begin{center}
\begin{tabular}{ll}
Save \texttt{key = 1} & \arraydiagram[1.5in]{11}\\[6pt]
Shift 5 to the right & \arraydiagram[1.5in]{12}\\[6pt]
Shift 4 to the right & \arraydiagram[1.5in]{13}\\[6pt]
Shift 2 to the right & \arraydiagram[1.5in]{14}\\[6pt]
Write 1 into the gap & \arraydiagram[1.5in]{10}
\end{tabular}
\end{center}
Repeated numbers during shifting are temporary: the key is safely stored in a variable. Gold tracks the active slot or gap.
\end{frame}

\begin{frame}{Insertion Sort: Describe the Procedure}
\small
\begin{enumerate}
 \item Start with the first element as a one-value sorted region.
 \item Save the next value in a variable called \texttt{key}.
 \item Walk left through the prefix while values are larger than \texttt{key}.
 \item Shift each larger value one position to the right.
 \item Write \texttt{key} into the gap and grow the sorted region.
\end{enumerate}
\begin{keyidea}Insertion sort makes room by shifting, rather than finding the overall minimum.\end{keyidea}
\end{frame}

\begin{frame}[fragile]{Insertion Sort: Python Implementation}
\small
\begin{block}{Sort the existing list in ascending order}
\footnotesize
\begin{verbatim}
def insertion_sort(values: list[int]) -> None:
    for position in range(1, len(values)):
        key = values[position]
        scan = position - 1
        while scan >= 0 and values[scan] > key:
            values[scan + 1] = values[scan]
            scan -= 1
        values[scan + 1] = key
\end{verbatim}
\end{block}
\begin{itemize}
 \item \texttt{position} identifies the next item to insert; \texttt{key} holds it safely.
 \item \texttt{scan} moves left through the sorted prefix.
 \item \texttt{and} tests \texttt{scan >= 0} first; when false, it skips \texttt{values[scan]}.
\end{itemize}
\end{frame}

\begin{frame}[fragile]{Trace the Code: Insertion Sort, First Key}
\small
\begin{columns}[T,onlytextwidth]
\begin{column}{.53\textwidth}
\begin{block}{Relevant code}
\scriptsize
\begin{verbatim}
position = 1
key = values[position]
scan = position - 1
while scan >= 0 and values[scan] > key:
    values[scan + 1] = values[scan]
    scan -= 1
values[scan + 1] = key
\end{verbatim}
\end{block}
\end{column}
\begin{column}{.44\textwidth}
\begin{block}{Values and assignments}
\footnotesize
\begin{tabular}{ll}
Initial & \texttt{[5,2,4,1,3]}\\
\texttt{key} & 2\\
\texttt{scan} & 0\\
Shift 5 & \texttt{[5,5,4,1,3]}\\
\texttt{scan -= 1} & $-1$\\
Insert 2 & \texttt{[2,5,4,1,3]}
\end{tabular}
\end{block}
\end{column}
\end{columns}
\begin{keyidea}The extra copy of 5 is temporary. When the loop ends, the saved key fills the gap.\end{keyidea}
\end{frame}

\begin{frame}{Trace the Code: Insertion Sort, A Larger Shift}
\small
Take \texttt{key = 1} at \texttt{position = 3} in \texttt{[2, 4, 5, 1, 3]}.
\begin{center}
\begin{tabular}{clll}
\texttt{scan} & Compare & Assignment & List afterward\\\hline
2 & $5>1$ & \texttt{values[3] = 5} & \texttt{[2,4,5,5,3]}\\
1 & $4>1$ & \texttt{values[2] = 4} & \texttt{[2,4,4,5,3]}\\
0 & $2>1$ & \texttt{values[1] = 2} & \texttt{[2,2,4,5,3]}\\
$-1$ & stop & \texttt{values[0] = 1} & \texttt{[1,2,4,5,3]}
\end{tabular}
\end{center}
\begin{keyidea}When \texttt{scan} reaches $-1$, the saved key belongs at index \texttt{0}.\end{keyidea}
\end{frame}

\begin{frame}{Trace the Code: Insertion Sort, Outer Loop}
\small
\begin{center}
\begin{tabular}{ccll}
\texttt{position} & \texttt{key} & Shifts & List after insertion\\\hline
1 & 2 & 5 & \texttt{[2,5,4,1,3]}\\
2 & 4 & 5 & \texttt{[2,4,5,1,3]}\\
3 & 1 & 5,4,2 & \texttt{[1,2,4,5,3]}\\
4 & 3 & 5,4 & \texttt{[1,2,3,4,5]}
\end{tabular}
\end{center}
\begin{keyidea}Each iteration inserts one new key into a prefix that was already ordered.\end{keyidea}
\end{frame}

\begin{frame}{Insertion Sort: Check Understanding}
\small
\begin{center}\arraydiagram[2in]{09}\end{center}
\begin{enumerate}
 \item Which three values are in the sorted prefix?
 \item What value is saved in \texttt{key}?
 \item Which existing values must shift right?
 \item At what index is \texttt{key} finally inserted?
\end{enumerate}
\begin{alertblock}{Remember}\texttt{key} is saved before any overwriting begins.\end{alertblock}
\end{frame}

\begin{frame}{Selection Sort vs. Insertion Sort}
\small
\begin{center}
\begin{tabular}{@{}p{.18\linewidth}p{.37\linewidth}p{.37\linewidth}@{}}
& \textbf{Selection sort} & \textbf{Insertion sort}\\[4pt]\hline
Next step & Find smallest value in unsorted suffix & Take next value and move it left into order\\[5pt]
Movement & One exchange after scanning & Shift larger values, then place key\\[5pt]
Sorted region & Contains smallest finished values & Contains values seen so far, in order\\[5pt]
Changes input? & Yes & Yes
\end{tabular}
\end{center}
\begin{keyidea}Both methods help us understand sorting. In everyday Python code, use built-in sorting.\end{keyidea}
\end{frame}

\sectionframe{Part 4: Python Sorts for You}

\begin{frame}[fragile]{The List Method: \texttt{.sort()}}
\small
\begin{block}{Modify the original list}
\footnotesize
\begin{verbatim}
values = [5, 2, 4, 1, 3]
values.sort()
print(values)  # [1, 2, 3, 4, 5]
\end{verbatim}
\end{block}
\begin{itemize}
 \item \texttt{.sort()} belongs to a list object.
 \item It rearranges the existing list in place.
 \item It returns \texttt{None}, not the sorted list.
\end{itemize}
\begin{alertblock}{Common mistake}\texttt{result = values.sort()} makes \texttt{result} equal to \texttt{None}.\end{alertblock}
\end{frame}

\begin{frame}[fragile]{The Function: \texttt{sorted()}}
\small
\begin{block}{Create a new sorted list}
\footnotesize
\begin{verbatim}
values = [5, 2, 4, 1, 3]
ordered = sorted(values)
print(ordered)  # [1, 2, 3, 4, 5]
print(values)   # [5, 2, 4, 1, 3]
\end{verbatim}
\end{block}
\begin{itemize}
 \item \texttt{sorted(values)} returns a new list; the input list remains unchanged.
 \item \texttt{sorted()} also accepts other iterables, such as a tuple or string.
 \item With objects, the new list still holds references to the original objects.
\end{itemize}
\begin{keyidea}Use \texttt{.sort()} to change the original; use \texttt{sorted()} for a separate list.\end{keyidea}
\end{frame}

\begin{frame}[fragile]{Descending Order: \texttt{reverse=True}}
\small
\begin{block}{Both sorting interfaces accept the option}
\footnotesize
\begin{verbatim}
values = [5, 2, 4, 1, 3]
values.sort(reverse=True)
print(values)  # [5, 4, 3, 2, 1]

ordered = sorted([5, 2, 4, 1, 3], reverse=True)
print(ordered)  # [5, 4, 3, 2, 1]
\end{verbatim}
\end{block}
\begin{itemize}
 \item \texttt{reverse=True} requests descending order.
 \item \texttt{values.reverse()} just flips the existing sequence; it does not sort it.
\end{itemize}
\end{frame}

\begin{frame}{Which Should I Use?}
\small
\begin{center}
\begin{tabular}{ll}
\textbf{Goal} & \textbf{Python expression}\\\hline
Reorder the existing list & \texttt{values.sort()}\\[6pt]
Make a new ordered list & \texttt{sorted(values)}\\[6pt]
Reorder existing list, descending & \texttt{values.sort(reverse=True)}\\[6pt]
Make a descending copy & \texttt{sorted(values, reverse=True)}
\end{tabular}
\end{center}
\begin{keyidea}You usually do not need to implement selection or insertion sort to sort real Python data.\end{keyidea}
\end{frame}

\sectionframe{Part 5: Sorting Objects with \texttt{lambda}}

\begin{frame}[fragile]{A List of Objects Needs an Ordering Rule}
\small
\begin{block}{The Student class from last lecture}
\footnotesize
\begin{verbatim}
class Student:
    def __init__(self, student_id, name, gpa):
        self.student_id = student_id
        self.name = name
        self.gpa = gpa
\end{verbatim}
\end{block}
\begin{itemize}
 \item A student has several attributes: ID, name, and GPA.
 \item Should students be ordered by ID? Alphabetically? By GPA?
 \item Python cannot infer which attribute we mean from this simple class.
\end{itemize}
\begin{keyidea}For object sorting, the programmer supplies the key used for comparisons.\end{keyidea}
\end{frame}

\begin{frame}[fragile]{The \texttt{key=} Argument}
\small
\begin{block}{Sort names by their length}
\footnotesize
\begin{verbatim}
names = ["Ada", "Christopher", "Bo"]
ordered = sorted(names, key=len)
print(ordered)  # ['Bo', 'Ada', 'Christopher']
\end{verbatim}
\end{block}
\begin{itemize}
 \item \texttt{key=} receives a function, not a precomputed result.
 \item Python calls that function once per item to obtain its sorting key.
 \item The return value determines the ordering: \texttt{len("Bo")} is 2.
\end{itemize}
\begin{keyidea}\texttt{key=} specifies what property to compare, not how to write a sorting loop.\end{keyidea}
\end{frame}

\begin{frame}[fragile]{What Is a \texttt{lambda}?}
\small
\begin{cpdefinition}A \texttt{lambda} expression creates a small function without giving it a name using \texttt{def}.\end{cpdefinition}
\begin{columns}[T,onlytextwidth]
\begin{column}{.47\textwidth}
\begin{block}{Named function}
\footnotesize
\begin{verbatim}
def get_gpa(student):
    return student.gpa
\end{verbatim}
\end{block}
\end{column}
\begin{column}{.49\textwidth}
\begin{block}{Equivalent lambda}
\footnotesize
\begin{verbatim}
lambda student: student.gpa
\end{verbatim}
\end{block}
\end{column}
\end{columns}
\begin{keyidea}A lambda receives input and evaluates one expression, whose result is returned automatically.\end{keyidea}
\end{frame}

\begin{frame}{Read \texttt{lambda student: student.gpa} Piece by Piece}
\small
\begin{center}
\texttt{lambda student: student.gpa}\medskip

\begin{tabular}{ll}
\textbf{Piece} & \textbf{Meaning}\\\hline
\texttt{lambda} & create a function value\\[6pt]
\texttt{student} & parameter: one input object\\[6pt]
\texttt{:} & separates input from the expression\\[6pt]
\texttt{student.gpa} & value returned for that object
\end{tabular}
\end{center}
\begin{block}{Example evaluation}Called with \texttt{Student(1042, "Luis", 3.8)}, this function returns \texttt{3.8}.\end{block}
\end{frame}

\begin{frame}[fragile]{Sort Student Objects by GPA}
\small
\begin{block}{The input list}
\footnotesize
\begin{verbatim}
students = [
    Student(1008, "Maya", 3.4),
    Student(1042, "Luis", 3.8),
    Student(1091, "Nora", 3.6),
]
\end{verbatim}
\end{block}
\begin{block}{Sort by the GPA attribute}
\footnotesize
\begin{verbatim}
by_gpa = sorted(students, key=lambda s: s.gpa)
\end{verbatim}
\end{block}
\begin{center}Resulting order: Maya (3.4) $\rightarrow$ Nora (3.6) $\rightarrow$ Luis (3.8)\end{center}
\end{frame}

\begin{frame}{How Does \texttt{key=lambda s: s.gpa} Work?}
\small
\begin{center}
\begin{tabular}{lcl}
\textbf{Object} & \textbf{Key returned} & \textbf{Sorted position}\\\hline
\texttt{Student(1008, "Maya", 3.4)} & 3.4 & first\\[6pt]
\texttt{Student(1042, "Luis", 3.8)} & 3.8 & third\\[6pt]
\texttt{Student(1091, "Nora", 3.6)} & 3.6 & second
\end{tabular}
\end{center}
\begin{itemize}
 \item The function receives each \texttt{Student} object as \texttt{s}.
 \item It returns that student's GPA as the comparison key.
 \item The result contains the original objects, not just GPAs.
\end{itemize}
\begin{keyidea}Changing the key changes the ordering rule without changing the objects.\end{keyidea}
\end{frame}

\begin{frame}[fragile]{Different Keys, Different Orders}
\small
\begin{block}{Same list of Student objects}
\footnotesize
\begin{verbatim}
# Alphabetical by name, new list
by_name = sorted(students, key=lambda s: s.name)

# Highest GPA first, original list changes
students.sort(key=lambda s: s.gpa, reverse=True)

# Increasing student ID, new list
by_id = sorted(students, key=lambda s: s.student_id)
\end{verbatim}
\end{block}
\begin{itemize}
 \item \texttt{key=} chooses the attribute; \texttt{reverse=True} reverses the order.
 \item \texttt{.sort()} versus \texttt{sorted()} determines whether the input changes.
 \item Equal keys keep their original relative order: Python sorting is \textbf{stable}.
\end{itemize}
\end{frame}

\begin{frame}{Connect \texttt{lambda} to Last Lecture}
\small
\begin{columns}[T,onlytextwidth]
\begin{column}{.48\textwidth}
\begin{block}{Last lecture: search predicate}\texttt{lambda s: s.gpa >= 3.8}\end{block}
Returns \texttt{True} or \texttt{False}: does this object match?
\end{column}
\begin{column}{.48\textwidth}
\begin{block}{Today: sorting key}\texttt{lambda s: s.gpa}\end{block}
Returns a GPA number: what value determines this object's order?
\end{column}
\end{columns}
\begin{keyidea}The lambda syntax is identical; the function's purpose and returned value differ.\end{keyidea}
\end{frame}

\begin{frame}{Common Mistakes with Object Sorting}
\small
\begin{enumerate}
 \item \texttt{key=lambda s: s.gpa()} is wrong if \texttt{gpa} is an attribute, not a method.
 \item \texttt{key=lambda s: s.gpa > 3.5} groups by a Boolean key; it does not order by GPA within those groups.
 \item \texttt{key=lambda s: s} asks Python to compare whole objects, which our basic \texttt{Student} class does not support.
 \item \texttt{students = students.sort(...)} replaces the variable with \texttt{None}.
\end{enumerate}
\begin{keyidea}A good key returns exactly the property you want the elements sorted by.\end{keyidea}
\end{frame}

\begin{frame}{Integrated Practice: Predict the Order}
\small
Suppose \texttt{students} contains Maya (3.4), Luis (3.8), and Nora (3.6).
\begin{enumerate}
 \item Who comes first after \texttt{sorted(students, key=lambda s: s.name)}?
 \item Who comes first after \texttt{sorted(students, key=lambda s: s.gpa)}?
 \item Who comes first with \texttt{key=lambda s: s.gpa, reverse=True}?
 \item Which changes the original list: \texttt{.sort()} or \texttt{sorted()}?
 \item Does the result contain \texttt{Student} objects or just numbers and names?
\end{enumerate}
\end{frame}

\begin{frame}{Takeaways}
\small
\begin{enumerate}
 \item Binary search relies on ordered data; sorting can provide that order.
 \item Both algorithms grow a sorted region from left to right.
 \item Selection sort finds the minimum remaining value and swaps it into place.
 \item Insertion sort saves the next key, shifts larger values, and inserts the key.
 \item \texttt{.sort()} changes the original list; \texttt{sorted()} returns a new list.
 \item \texttt{reverse=True} requests descending order.
 \item \texttt{lambda} creates a short function; \texttt{key=lambda s: s.gpa} orders by GPA.
\end{enumerate}
\begin{keyidea}Understand the algorithms; use built-in sorting in practice.\end{keyidea}
\end{frame}
\end{document}
